\section{The explicit power-cover lemma}

\subsection{Unit-ideal covers}

For a ring $A$, a finite family of principal opens $D(f_0),\ldots,D(f_{r-1})$ covers
$\Spec(A)$ exactly when the elements $f_i$ generate the unit ideal.  In the present framework the
cover is not represented by quantifying over prime ideals.  It is represented by coefficients
$c_i\in A$ and an equality witness for
\begin{equation}\label{eq:unit-cover}
                 \sum_{i<r} c_i f_i=1.
\end{equation}
This certificate is finite, stable under ring maps, and directly usable by an arithmetic proof.

Gluing fractions usually requires a common denominator power.  The same principal opens are also a
cover after replacing every $f_i$ by $f_i^N$, but an iPRA proof needs explicit coefficients.

\begin{theorem}[Power-cover constructor]\label{thm:power-cover}
Let $r\geq1$, let $f_i,c_i\in A$ satisfy \eqref{eq:unit-cover}, and let $N\geq1$.  There is an
explicit primitive-recursive construction of coefficients $d_i$ such that
\begin{equation}\label{eq:power-cover}
                  \sum_{i<r}d_i f_i^N=1.
\end{equation}
The construction is uniform in the finite ring codes, $r$, and $N$.
\end{theorem}

\begin{proof}
Set
\[
                    M=r(N-1)+1.
\]
Expand the finite power of \eqref{eq:unit-cover}:
\begin{equation}\label{eq:multinomial-cover}
  1=\left(\sum_{i<r}c_i f_i\right)^M
   =\sum_{\substack{\alpha\in\N^r\\ |\alpha|=M}}
      \binom{M}{\alpha}\prod_{j<r}c_j^{\alpha_j}f_j^{\alpha_j}.
\end{equation}
For every exponent vector $\alpha$ of total degree $M$, some coordinate satisfies
$\alpha_i\geq N$.  Otherwise
$|\alpha|\leq r(N-1)<M$.  Choose the least such $i$ and assign the corresponding term of
\eqref{eq:multinomial-cover} to bucket $i$.  Factor $f_i^N$ from every term in bucket $i$ and let
$d_i$ be the sum of the residual factors.  Then the buckets partition the finite expansion, so
\eqref{eq:power-cover} follows.
\end{proof}

The proof is constructive even over rings with zero divisors.  It uses no radical membership,
prime ideal, compactness, or excluded middle beyond comparisons of natural-number exponents.

\subsection{Exact coefficient formula}

For reproducibility, define
\[
  \iota(\alpha)=\min\set{i<r:\alpha_i\geq N}.
\]
Then
\begin{equation}\label{eq:di-formula}
 d_i=\sum_{\substack{|\alpha|=M\\ \iota(\alpha)=i}}
 \binom{M}{\alpha}
 c_i^{\alpha_i}f_i^{\alpha_i-N}
 \prod_{j\neq i}c_j^{\alpha_j}f_j^{\alpha_j}.
\end{equation}
The finite weak compositions of $M$ into $r$ parts can be enumerated by nested primitive recursion.
The multinomial coefficient is computed from bounded products and exact divisions, or avoided
entirely by an iterative expansion that carries multiplicities by repeated addition.

\subsection{An iPRA-friendly iterative proof}

Formula \eqref{eq:di-formula} is mathematically transparent, but the object proof need not formalize
factorials.  A more direct implementation maintains a finite polynomial table indexed by exponent
vectors with coordinates at most $M$.  Starting from the constant table $1$, perform $M$ rounds of
multiplication by $\sum c_i f_i$.  The quantifier-free invariant states:

\begin{enumerate}[label=(\roman*)]
\item every occupied exponent vector has total degree equal to the round number;
\item the sum of table entries times their monomials equals the current power;
\item each entry is the accumulated coefficient obtained by the fixed update rule.
\end{enumerate}

After round $M$, scan the finite table, choose the least coordinate at least $N$, and add the factored
term to the corresponding $d_i$.  A second quantifier-free invariant states that the processed prefix
of the table equals the partial sum $\sum_i d_i f_i^N$.  This route uses only finite arrays, bounded
comparison, and ring operations already available in the presentation.

\subsection{Examples and validation}

For $f_0=2$, $f_1=3$, $c_0=-1$, and $c_1=1$, the basic cover certificate is
$-2+3=1$.  At power $N=2$, the constructor used in the executable companion returns
\[
                   d_0=7,\qquad d_1=-3,
\]
and indeed
\[
                    7\cdot2^2-3\cdot3^2=28-27=1.
\]

The same algorithm was tested in $\Z/12\Z$ with $f_0=3$, $f_1=4$ and
$c_0=3$, $c_1=1$, since $3\cdot3+1\cdot4=1\pmod {12}$.  This test is useful because neither
principal element is cancellable and localization witnesses genuinely need annihilating powers.

\subsection{Functoriality}

If $\phi:A\to B$ is a certified ring map, applying $\phi$ to the input cover and to every operation
in \eqref{eq:di-formula} gives exactly the power-cover coefficients constructed in $B$.  Thus the
constructor is functorial up to the supplied equality witnesses.  This fact later makes atlas
refinement and base change finite certificate transformations rather than new existence proofs.
